\documentclass[11pt,a4paper]{article}
\usepackage[T1]{fontenc}
\usepackage{lmodern,amsmath,amssymb,amsthm,booktabs,microtype,hyperref}
\usepackage[margin=27mm]{geometry}
\setlength{\emergencystretch}{2em}
\newcommand{\Q}{\mathbf Q}
\newcommand{\Z}{\mathbf Z}
\newcommand{\Tr}{\operatorname{Tr}}
\newcommand{\N}{\operatorname{N}}
\newtheorem{theorem}{Theorem}
\newtheorem{lemma}[theorem]{Lemma}
\newtheorem{proposition}[theorem]{Proposition}
\title{Unrestricted auxiliary points in the derivative detector}
\author{}\date{17 November 2026}
\begin{document}\maketitle
\begin{abstract}
The auxiliary points in the derivative detector may range over the
entire elliptic curve, subject to its affine guards, while the root
point remains in the index subgroup. This preserves soundness and
the complete eleven-variable construction. Cubic auxiliary point
charts give an integer circuit of degree at most $232094337$, with
integer content one. The same enlargement applies to the sound
nine-variable actual-trace source; its global rational realization
remains a separate requirement.
\end{abstract}

\section{Auxiliary point charts}
Retain the arithmetic premises, coefficient shift and notation of
\cite{Derivative,Criterion}. In particular,
\[
 K=\Q(\sqrt2),\qquad D=11-6\sqrt2,\qquad
 E:y^2=x^3-Dx,\qquad 2E(K)=\Z P,
\]
\[
 P=\left(\frac{57-28\sqrt2}{8},
          \frac{-660+417\sqrt2}{32}\right).
\]
The root point $A=[n]P$ retains the square-abscissa chart
$x=g^2$, $y=gk$, $k^2=g^4-D$. For each auxiliary point $H,U$,
retain two rational coordinates $h_0,h_1$ and put
\begin{equation}\label{eq:chart}
 h=h_0+h_1\sqrt2,\qquad x=h,\qquad y=j,\qquad
 j^2=h^3-Dh=a+b\sqrt2,
\end{equation}
where
\begin{align*}
 a&=h_0^3+6h_0h_1^2-11h_0+12h_1,\\
 b&=3h_0^2h_1+2h_1^3+6h_0-11h_1.
\end{align*}
Retain the nonzero point ordinates and abscissas, the affine
difference denominators, and all original detector denominators.
Thus points at infinity and points of order two remain outside
these affine charts.

Use the free rank-four algebra of \cite{Unit}, with basis
$1,\mathsf U,\mathsf V,\mathsf W$ and multiplication table
\[
\begin{gathered}
 \mathsf U^2=2a+2\mathsf W,\qquad
 \mathsf V^2=a-\mathsf W,\qquad \mathsf U\mathsf V=2b,\\
 \mathsf U\mathsf W=a\mathsf U-2b\mathsf V,\qquad
 \mathsf V\mathsf W=b\mathsf U-a\mathsf V,\qquad
 \mathsf W^2=a^2-2b^2.
\end{gathered}
\]
Then $j=(\mathsf U+\mathsf V\sqrt2)/2$ and
$\N_{K/\Q}(j)=\mathsf W$. The existing inversion of the ordinate
therefore retains $\mathsf W\ne0$; it forces $a^2-2b^2\ne0$.

\begin{lemma}\label{lem:charts}
On these guards, \eqref{eq:chart} has exactly the guarded affine
points of $E(K)$ as its rational characters. Its point algebra is
etale of rank four. The original nonzero points of $2E(K)$ embed
with a bijection of all geometric characters.
\end{lemma}
\begin{proof}
After adjoining $\sqrt2$, the algebra is the tensor product of the
two square equations with radicands $a+b\sqrt2$ and
$a-b\sqrt2$. Their product is nonzero on the ordinate guard.
The character description in \cite{Unit} applies without changing
the multiplication table. Conversely, every guarded affine point
has $h=x$ and $j=y$, so is represented.

For an original chart $x=g^2$, $y=gk$, set $h=g^2$ and $j=gk$.
The guard $x\ne0$ makes $g$ invertible in both field slots.
Consequently multiplication by $g$ identifies the old square-root
algebra $k^2=g^4-D$ with the new one $j^2=g^6-Dg^2$ at every
geometric character. All point values, difference values and
detector values are unchanged under this identification.
\end{proof}

For the retained primitive element
$z=(2a+b)\mathsf U-2(a+b)\mathsf V$, put
\[
 a'=6a^3-9ab^2-4b^3,\qquad
 b'=-4a^3+6ab^2+3b^3.
\]
Its quartic is $Z^4-4a'Z^2+8(b')^2$, with discriminant
$2^{15}(b')^2(2a^2-b^2)^4(a^2-2b^2)^2$.
At rational retained coordinates, $a^2-2b^2\ne0$ implies
$2a^2-b^2\ne0$. Also $b'\ne0$: for $b\ne0$ use the
Eisenstein polynomial $4T^3-6T-3$, and for $b=0$ use $a\ne0$.
Hence the same two-stage point separator has nonzero radicands.
No extra retained variable is needed.

\section{The local arithmetic implication}
Keep the six tests
\[
\begin{gathered}
 C_1=f(H),\quad C_2=f(U),\quad
 C_3=\overline{f(H-P)},\quad C_4=\overline{f(U-A)},\\
 C_5=r-r',\quad C_6=r-X,\qquad
 r=\frac{\zeta(U)}{\zeta(H)},\qquad
 r'=\overline{\frac{\zeta(U-A)}{\zeta(H-P)}},
\end{gathered}
\]
where $f=1/x$ and $\zeta=-x/y$.

\begin{lemma}\label{lem:formal}
The local extraction implication of the derivative detector holds
with arbitrary guarded $H,U\in E(K)$ and $A=[n]P$.
\end{lemma}
\begin{proof}
There exist $h',j'\in\Z$ such that
$2H=[h']P$ and $2U=[j']P$. Put $h=h'/2$ and $j=j'/2$.
At every odd prime, $h,j,h-1,j-n$ are locally integral, and one
of $h,h-1$ is a unit. The difference guards exclude $h=0,1$
and $j=0,n$: any corresponding zero doubled point would be a
point of order two, whose ordinate is guarded nonzero.

Extend the local elliptic logarithm to the elliptic group by
first multiplying into the formal group and then dividing by
that integer. This homomorphism kills torsion. Thus at each
relevant embedding its exact ratios are
\[
 \frac{\log_E(U)}{\log_E(H)}=\frac jh,
 \qquad
 \frac{\log_E(U-A)}{\log_E(H-P)}=\frac{j-n}{h-1}.
\]
The simultaneous positivity of the first four tests puts the
appropriate points in the respective formal groups. At the good
odd primes retained by the source, the existing comparison
$\log_E(Q)=\zeta(Q)(1+O(p^4))$ therefore gives the same truncated
ratios as before. Because $h$ or $h-1$ is a local unit, one ratio
is integral; $v_p(r-r')>0$ makes the other integral as well.
The positive target test gives $X\equiv r\equiv r'\pmod p$.
Multiplying by $h$ and $h-1$ yields
\[
 Xh\equiv j,\qquad X(h-1)\equiv j-n\pmod p.
\]
Subtracting proves $X\equiv n\pmod p$. The only additional
index denominator is two, already among the exceptional primes.
\end{proof}

\begin{theorem}\label{thm:transfer}
Under ordinary rational $abc$ and the retained arithmetic premises
of \cite{Derivative}, the full auxiliary eleven-variable source
defines $\Z$ for every sufficiently large fixed odd shift $B$.
Its integer circuit has degree at most $232094337$ and content
one. The actual-trace replacement of \cite{Criterion} remains a
sound source with nine variables and has a larger auxiliary
witness space. Its completeness still requires global rational
realization of the actual square equation.
\end{theorem}
\begin{proof}
The valuation argument extracting positivity of the six tests
depends only on their local values and is unchanged. Lemma
\ref{lem:formal} supplies exactly the former local divisibility
statement. The scale, derivative weight, root point $A$, and its
index-height inequality are unchanged. Hence the same height
absorption proves soundness, for the same kind of sufficiently
large fixed odd shift.

For the complete eleven-variable source, choose an original
prepared witness. Lemma \ref{lem:charts} maps all of its point
characters bijectively to the new chart and preserves every
old guard. The newly necessary auxiliary radicand guards hold
because $j=gk$ is nonzero at every old point character. The pole
parameters and terminal root are unchanged. Thus the old whole-
character completeness proof supplies new complete witnesses.
The finite-algebra encoder retains the same ranks and variables.

The nine-variable square test has the same local implication
and the same finite-algebra terminal-root allowance. Consequently
its soundness transfers. The above eleven-variable completeness
argument uses the intrinsic pole parameters, so does not supply
the missing rational square at the actual trace.
\end{proof}

\section{Integer circuit bounds}
Assign point basis weights $(0,2,2,4)$ to the root block and
$(0,2,2,3)$ to each auxiliary block. Every relation in the displayed
table respects these weights, since the new $a,b$ have degree three.
For coefficient operator norms use diagonal basis scales
$(1,16,16,80)$ in all three blocks. The coefficient norms
$\|a\|_1=30$, $\|b\|_1=22$ bound each new block. The old root
block has the smaller values $28,18$. Direct inspection of the
six table relations bounds multiplication by
$\mathsf U,\mathsf V,\mathsf W$ by $16,16,80$, respectively.
For example, the final column for $\mathsf U$ is bounded by
$(30+44)16/80<16$, and its second column by
$(60+160)/16<16$.

The point norm and primitive radicands satisfy the coefficient
majorants
\[
 30^2+2\cdot22^2=1868<2^{11},
\]
\[
 2(6\cdot30^3+9\cdot30\cdot22^2+4\cdot22^3)
 +160(2\cdot30^2+22^2)=1035984<2^{20}.
\]
The primitive root has operator norm at most
$(4\cdot30+3\cdot22)16=2976<2^{12}$.
The auxiliary quadratic stages have degrees six and nine, with
root weights three and five. The original root stages retain
degrees eight and twelve, with root weights four and six.
The terminal root may vanish; all preceding stages remain units.

Applying the subalgebra norm and terminal-root specifications of
\cite{Subalgebra,Terminal} gives
\[
 \deg G_*\le1526992,\qquad \deg\Omega\le286244,
\]
and therefore
\[
 \deg\bigl(X\det(I-E_0G_*M_\Omega)\bigr)
 \le1+128(1526992+286244+1)=232094337.
\]
The encoder has degree at most $128$, and setting $E_0=0$ gives
exactly $X$. Integer content is one and the support bound remains
strictly less than $2^{264}$. For the algebraic fixture $B=3$, the
coefficient norm bound is $2^{487256960}$. This fixture is not
asserted to be a sufficient defining shift.

The accompanying compiler records all $1990$ integer arithmetic
nodes, $32$ native guards and $38$ retained norm factors. An
independent iterative traversal checks each factor, subalgebra,
degree, coefficient bound and denominator. Direct chord-and-tangent
evaluation checks all $64$ point characters for each of three
actual tuples at two primes. The auxiliary indices are
$1/2,3/2,5/2$, so these tests include points unavailable to the
former auxiliary chart. These finite checks verify the circuit;
they do not establish the retained external arithmetic inputs
or global nine-variable realization.

\begin{thebibliography}{9}\raggedright
\bibitem{Derivative} \emph{An eleven-variable construction from a derivative detector},
project source, \nolinkurl{papers/conditional_few_variables/derivative_detector11_20261117_v63.tex}.
\bibitem{Criterion} \emph{One actual-trace square in a derivative detector},
project source, \nolinkurl{papers/pole_parameters/single_trace_derivative_criterion_20261117_v64.tex}.
\bibitem{Unit} \emph{A single norm determinant for the 25- and 13-variable constructions},
project source, \nolinkurl{papers/finite_algebra/unit_norm_20261114_v28.tex}.
\bibitem{Subalgebra} \emph{Subalgebra norm guards and integral quadratic separators},
project source, \nolinkurl{papers/finite_algebra/subalgebra_norms_20261114_v30.tex}.
\bibitem{Terminal} \emph{A possibly vanishing terminal quadratic root},
project source, \nolinkurl{papers/finite_algebra/terminal_quadratic_root_20261115_v40.tex}.
\end{thebibliography}
\end{document}
